\chapter{Marco teórico y antecedentes}\label{cap:2}

\pendiente{capítulo nuevo. Abajo, el esquema propuesto con lo que debe cubrir cada
sección. Las referencias marcadas ``verificar'' son sugerencias que aún
no están en el paper.}

\section{Precipitación diaria, días secos y rachas secas}\label{sec:mt_rachas}

\begin{itemize}
\tightlist
\item
  Definición de día húmedo y seco, y la convención de 1 mm (ETCCDI; Zhang et al., 2011).
\item
  Índices escalares de persistencia (CDD, número de días secos, duración media) y sus limitaciones (Sillmann et al., 2013a, 2013b; Giorgi et al., 2011, 2019).
\item
  Distribuciones de duración: colas pesadas y desviaciones de un proceso sin memoria (Zolina et al., 2013; Lana et al., 2006; Herrera-Marín et al., 2026; Stechmann y Neelin, 2014; Martinez-Villalobos y Neelin, 2019).
\end{itemize}

\section{Clima de Chile y su gradiente norte--sur}\label{sec:mt_clima}

\begin{itemize}
\tightlist
\item
  Controles del clima: anticiclón del Pacífico Sur, vientos del oeste, los Andes y la estacionalidad de la precipitación (verificar: Garreaud, 2009).
\item
  Variabilidad interanual y decadal: ENSO, PDO y SAM (verificar: Montecinos y Aceituno, 2003).
\item
  La megasequía y el secado reciente (Garreaud et al., 2017, 2020; Boisier et al., 2016, 2018) y las tendencias de dry spells (Latoja et al., 2025).
\end{itemize}

\section{Modelos climáticos, escenarios y sesgos}\label{sec:mt_modelos}

\begin{itemize}
\tightlist
\item
  GCM y RCM; reducción de escala dinámica frente a estadística; CMIP6 y CORDEX.
\item
  Escenarios SSP: diferencia entre SSP5-8.5 y SSP3-7.0, incluidos los aerosoles (verificar: O'Neill et al., 2016; Riahi et al., 2017).
\item
  Sesgo de llovizna de los modelos y sus efectos en las rachas secas (Stephens et al., 2010; Sun et al., 2006; Chen et al., 2021; Martinez-Villalobos et al., 2022; Petrova et al., 2024).
\item
  Proyecciones de precipitación y sequía para Chile (Araya-Osses et al., 2020; Bozkurt et al., 2018; Salazar et al., 2024).
\end{itemize}

\section{Área de estudio}\label{sec:area}

\pendiente{describir las cuatro ventanas y las cuatro regiones con su precipitación
anual y su estacionalidad (CR2MET 1980--2014: Coquimbo 247 mm, O'Higgins
956 mm, La Araucanía 2 038 mm, Los Lagos 2 734 mm; entre 85 \% y 97 \%
de la lluvia cae entre marzo y noviembre). Agregar un mapa de
precipitación anual de CR2MET.}

\section{Herramientas estadísticas}\label{sec:mt_estadistica}

\begin{itemize}
\tightlist
\item
  PDFs empíricas y bins logarítmicos para distribuciones de cola pesada (verificar: Newman, 2005).
\item
  Bootstrap y remuestreo por bloques (años) para datos dependientes (verificar: Efron y Tibshirani, 1993; Wilks, 2011).
\item
  Risk ratio como razón de conteos de excedencias y su interpretación.
\end{itemize}
